\section{\sys{} v2: An Economy for Agents}
\label{sec:system}

\sys{} v2 rebuilds v1~\cite{pai2025hashiru} on open protocols. A user's request goes to a CEO agent, which may hire
\emph{managers} (agents that use tools and may hire) and \emph{workers} (agents that answer once), to a configurable
depth (two levels below the CEO by default). Every agent is its own A2A service with its own Agent Card, every tool is
served over MCP, and \sys{} itself is reachable as an OpenAI-compatible model, an A2A agent and an MCP server. One
process serves everything on one port. Every model call, tool call and message is written to a JSONL trace.

\subsection{From v1 to v2}
v1 had one level of hierarchy, a hand-written table of resource units, and no analysis of the local-against-cloud
trade (its own stated limitations). v2 lets managers hire with money they were given, prices every cloud call from
LiteLLM's price map~\cite{berriai2026litellm} (local calls cost no dollars but hold the memory their model really uses,
measured from Ollama~\cite{ollama2026ollama}), and is measured by \bench{}. Rules that v1 left to prompts (who may
spend what, where private data may go) are enforced in code: in our experience with an earlier prototype, rules stated
only in prompts were broken.

\subsection{Pools, reservations and refunds}
\label{sec:pools}
Each agent $a$ holds a reservation $R_a$ carved from its hirer's pool. What it can still spend is
\begin{equation}
\mathrm{left}(a) = R_a - \mathrm{spent}(a) - \sum_{c \in \mathrm{hires}(a)} \mathrm{hold}(c),
\label{eq:left}
\end{equation}
where a live hire holds its whole reservation and a fired one holds only what it (and everyone below it) actually
spent. Firing therefore returns the unspent part, and two hires running at once cannot spend the same dollars: each
holds its reservation from the moment it is hired. Table~\ref{tab:pool} works one session through.

\begin{table}[t]
\centering
\caption{One session's pools (dollars). The CEO starts with \$0.50.}
\label{tab:pool}
\footnotesize
\begin{tabular}{@{}lrr@{}}
\toprule
Step & CEO left & m1 left \\
\midrule
Session opens, \texttt{budget\_usd} = 0.50 & 0.50 & \\
CEO spends 0.01 thinking & 0.49 & \\
CEO hires manager m1 (0.20) and worker w1 (0.10) & 0.19 & 0.20 \\
m1 hires w2 (0.05) and spends 0.02 & 0.19 & 0.13 \\
w1 spends 0.03, then the CEO fires it & 0.26 & 0.13 \\
\bottomrule
\end{tabular}
\end{table}

The checks are in code, not in prompts: a hire cannot reserve more than its hirer has left, an agent whose pool is
spent makes no further cloud call, and a remote agent is paid from its hirer's reservation, which must cover its
declared price before each message. A local hire is refused if its model does not fit in free memory (or in a cap,
\texttt{memory\_mb}). Separately, \texttt{max\_usd} caps real spend for the whole process. Because every reservation,
call and fire is in the trace, the ledger can be rebuilt from the trace alone, so budget rules can be checked after the
fact.

\subsection{Hiring by declared difficulty}
A hiring agent may name a model, or declare how hard a job is (easy, medium or hard) and let the hub choose. The hub
turns a difficulty into a list of models, best first, and hires the first that the modes, the hirer's pool and the
data's sensitivity allow: an easy job tries small local models, then cheap cloud models; a hard job the dearest cloud
models first. A worker that fails or answers empty moves up its list and gets the message again, up to a limit
(\texttt{max\_escalations}, two by default). Which models exist, and their tiers, come from the machine itself:
Ollama's pulled models that call tools, and every cloud model in LiteLLM's map whose provider key is set, tiered by
price or by size. Declared difficulty is the agent's own judgment, so \bench{} measures it against the hidden
difficulty of each job (\S\ref{sec:evaluation}).

\subsection{Make or buy over open protocols}
Besides models, a manager may hire a remote A2A agent as a worker (\texttt{a2a/<name>}). The hub reads the agent's
card, prices it through a small A2A extension (a card declares a price per message or per token; a reply reports what
the request cost), charges the larger of the declared price and the reported cost to the hire's pool, and pins the
agent to its configured URL. Since a \sys{} hub is itself an A2A agent, one hub can hire another: a whole team, with
its own CEO and pools, appears to the hirer as one priced worker. This turns hiring into make or buy: an agent can make
an answer with a model it pays per token, or buy it from an outside agent at a price.

\subsection{Quotes, authority and sensitivity}
Three further mechanisms read naturally in economic terms. \emph{Plan mode is a quote:} the CEO may only read and
propose; the hub previews each proposed hire as it would be made (the model the router picks, what it reserves, the
memory it takes) and checks the totals against the pool, and approving the plan approves its hires. An approved plan's
steps compile to a graph (LangGraph~\cite{langchain2026langgraph}) that hires, runs each step once its inputs exist,
fires each hire after its last step, and lets the CEO answer once. \emph{Approvals are spending authority:} a hire or a
turn above a threshold waits for the user. \emph{Sensitivity limits where work may go:} private data never reaches a
cloud or remote model as it is; it is masked with placeholders on the way out and restored on the way back (Presidio's
recognizers~\cite{presidio2026presidio} and a local model find the values), or refused.

\subsection{Implementation}
\sys{} v2 is 6,323 lines of Python 3.13 with exact pins for its protocol and model libraries, and 446 tests that run
without API keys or Ollama. The two changes \bench{} needed (a switch that takes the web tools away, so that no answer
can be looked up, and the memory cap) were added test-first. Running \bench{}'s first real calls found one defect:
a reasoning model (\texttt{gpt-6-luna}) rejects temperature 0, so \sys{} now lets LiteLLM drop that parameter for such
models only.
